\documentclass[11pt]{article}
\usepackage{amsmath,amssymb,hyperref}
\usepackage{geometry}
\usepackage{booktabs}
\usepackage{url}
\Urlmuskip=0mu plus 1mu
\geometry{margin=1in}

% Write-up for Experiments/SU3_Adjoint_SoV_ScalarProduct_L2.wb.
% Twisted continuation: SU3_Adjoint_Twisted_SoV_L2.tex.

\newcommand{\Op}{\mathcal{O}}
\newcommand{\sh}[1]{^{[#1]}}
\newcommand{\gl}{\mathfrak{gl}}
\newcommand{\br}[1]{\langle #1\rangle}

\title{Functional separation of variables at zero twist\\
for the $\gl(3)$ spin chain with two adjoint sites}
\author{Notes accompanying \texttt{Experiments/SU3\_Adjoint\_SoV\_ScalarProduct\_L2.wb}}
\date{14 September 2026}

\begin{document}
\maketitle

\begin{abstract}
We construct the functional SoV scalar product for the untwisted rational $\gl(3)$ spin chain of
length two with adjoint representations $[2,1,0]$ at both sites. This is the first
non-rectangular case. The transfer matrices are built explicitly on the $64$-dimensional space.
Their highest-weight eigenvalues depend on three independent differences of integrals of motion,
compared with two for the rectangular $\gl(4)$ chain. The orthogonality relation therefore gives a
$3\times3$ matrix. Its determinant vanishes for every pair of distinct states, including the two
octets, which share their weight, their Casimirs and $T_1$, and it is non-zero for each state.
With every row integrated over the full contour, the vacuum-normalised norm depends on the
inhomogeneities, unlike in $\gl(4)$. We compare it with the bilinear norm of the explicit
eigenvectors. We show that it cannot be converted into that norm by values of the gauge-invariant
$Q$-functions at pole-ladder points, and we record a bounded search that confirms this.
\end{abstract}

\tableofcontents

\section{The chain}

Each site carries the $8$-dimensional $\gl(3)$ irrep $\lambda=[2,1,0]$, realised in the
Gelfand--Tsetlin basis. The generators are taken over unchanged from \texttt{SU4\_L1.wb}; that
construction works for any rank. With inhomogeneities $\theta_1,\theta_2$ and shift $h$ we set
$f\sh{\pm2n}(u)=f(u\pm nh)$ and
\begin{equation}
L_{ij}(u)=u\,\delta_{ij}-h\,E_{ji},\qquad
T(u)=L^{(1)}(u-\theta_1)\,L^{(2)}(u-\theta_2).
\end{equation}
The transfer matrices are sums of principal quantum minors,
\begin{equation}
T_{a,1}(u)=\sum_{|I|=a}\ \sum_{\sigma\in S_a}\operatorname{sgn}\sigma\,
T_{I_{\sigma1}I_1}(u)\,T_{I_{\sigma2}I_2}(u-h)\cdots T_{I_{\sigma a}I_a}(u-(a-1)h).
\end{equation}
With
\begin{equation}
\nu_i(u)=(u-\theta_1-h\lambda_i)(u-\theta_2-h\lambda_i),\qquad \lambda=(2,1,0),
\end{equation}
we verified the following exactly and symbolically in $u,\theta_1,\theta_2,h$:
$[T_{a,1}(u),T_{b,1}(v)]=0$, the quantum determinant
$T_{3,1}=\nu_1\sh{-4}\nu_2\sh{-2}\nu_3$ is central, and every $T_{a,1}$ commutes with the total
$\gl(3)$ generators.

\section{Highest-weight states and integrals of motion}

The tensor product decomposes as $8\otimes8=27\oplus10\oplus\overline{10}\oplus8\oplus8\oplus1$,
with highest weights $\{4,2,0\}$, $\{3,3,0\}$, $\{4,1,1\}$, $\{3,2,1\}^{\times2}$ and $\{2,2,2\}$.
At zero twist the transfer matrices act on each highest-weight space. Only the $\{3,2,1\}$ space
is two-dimensional.

\paragraph{Operator identities.} As $64\times64$ matrices,
\begin{align}
[u^4]T_2&=3,\qquad [u^3]T_2=-6(3h+\theta_1+\theta_2),\\
[u^2]T_2&=[u^0]T_1+3\big(10h^2+\theta_1^2+3\theta_1\theta_2+\theta_2^2+8h(\theta_1+\theta_2)\big),
\end{align}
while $[u^1]T_2$ and $[u^0]T_2$ do not lie in the span of $[u^0]T_1$ and the identity.
The subleading coefficient of $T_1$ depends only on $\sum\mu_i=6$. Hence, for any two states,
\begin{equation}
\Delta T_1=\Delta t_1,\qquad \Delta T_2=\Delta t_1\,u^2+\Delta t_{21}\,u+\Delta t_{20},
\label{eq:dT}
\end{equation}
with three independent constants $(\Delta t_1,\Delta t_{21},\Delta t_{20})$. $T_1$ does not
distinguish $10$ from $\overline{10}$, nor the two octets. On the octet space, $T_2$ has eigenvalues
differing by the constant $3h^3\sqrt{4h^2+5(\theta_1-\theta_2)^2}$.

\paragraph{Parameter point.} We take $\theta_1=\tfrac13$, $\theta_2=-\tfrac8{57}$, $h=1$. Then
$\theta_1-\theta_2=\tfrac9{19}\notin\mathbb Z$ and the square root equals $\tfrac{43}{19}$, so all
eigenvalues and $Q$-functions are rational. The rational points of $4+5t^2=s^2$ are
$t=(m-1)(m-5)/(m^2-5)$; we also use $t=-\tfrac7{11}$ and $t=-\tfrac{33}{19}$.

\section{Baxter equations and $Q$-functions}

The finite-difference operators are
\begin{align}
\Op&=1-\frac{T_{1}}{\nu_1}\mathcal D^{-1}+\frac{T_{2}}{\nu_1\nu_1\sh{-2}}\mathcal D^{-2}
-\frac{T_{3}}{\nu_1\nu_1\sh{-2}\nu_1\sh{-4}}\mathcal D^{-3},\\
\Op^\dagger&=1-\frac{T_{1}\sh{2}}{\nu_3\sh{2}}\mathcal D+\frac{T_{2}\sh{4}}{\nu_3\sh{2}\nu_3\sh{4}}\mathcal D^{2}
-\frac{T_{3}\sh{6}}{\nu_3\sh{2}\nu_3\sh{4}\nu_3\sh{6}}\mathcal D^{3},
\end{align}
and we solve $\Op\,q_i\sh{2}=0$ and $\Op^\dagger q^j=0$. Once the eigenvalues are known these
equations are linear in $q$. Each has a three-dimensional polynomial kernel. We fix the gauge by
row reduction from the top degree, so that each solution is monic and has zero coefficient at the
degrees of the lower solutions. The degrees are
\begin{equation}
\deg q_i=(4-\mu_i)+(i-1),\qquad \deg q^j=\mu_j+(3-j),
\end{equation}
as listed in Table~\ref{tab:deg}.
\begin{table}[h]
\centering
\begin{tabular}{lcc}
\toprule
state & $\deg q_1,q_2,q_3$ & $\deg q^1,q^2,q^3$\\
\midrule
$\{4,2,0\}$ & $0,3,6$ & $6,3,0$\\
$\{3,3,0\}$ & $1,2,6$ & $5,4,0$\\
$\{4,1,1\}$ & $0,4,5$ & $6,2,1$\\
$\{3,2,1\}$ (both) & $1,3,5$ & $5,3,1$\\
$\{2,2,2\}$ & $2,3,4$ & $4,3,2$\\
\bottomrule
\end{tabular}
\caption{Degrees of the gauge-fixed direct and dual Baxter polynomials.}
\label{tab:deg}
\end{table}

\section{Measure and the orthogonality relation}

For $\lambda=(2,1,0)$,
\begin{equation}
\frac{\Gamma[\nu_1\sh{2}]}{\Gamma[\nu_3\sh{2}]}\ \propto\ \mu(u)=\prod_{\alpha=1}^2\frac{1}{(u-\theta_\alpha)(u-\theta_\alpha-h)},
\qquad \frac{\mu\sh2}{\mu}=\frac{\nu_1\sh2}{\nu_3\sh2}.
\end{equation}
The second relation is what makes $\oint\mu\,g\,\Op f=\oint\mu\,f\,\Op^\dagger g$ hold on a contour
that encloses all poles. The identity and the quantum determinant cancel in $\Op^A-\Op^B$. With
$E_2=\mu/(\nu_1\nu_1\sh{-2})=\prod_\alpha\prod_{n=0}^{3}(u-\theta_\alpha-nh)^{-1}$ and
\eqref{eq:dT}, this gives
\begin{equation}
0=\oint E_2\,q^j_A\big(-\Delta T_1\,\nu_1\sh{-2}q_i^B+\Delta T_2\,q_i^{B[-2]}\big)
=R^{AB}_{i,j}\cdot(\Delta t_1,\Delta t_{21},\Delta t_{20})^{T},
\end{equation}
where
\begin{equation}
R^{AB}_{i,j}=\Big(\oint E_2\,q^j_A\big(u^2q_i^{B[-2]}-\nu_1\sh{-2}q_i^B\big),\
\oint E_2\,q^j_A\,u\,q_i^{B[-2]},\ \oint E_2\,q^j_A\,q_i^{B[-2]}\Big).
\end{equation}
Since $\theta_1-\theta_2\notin h\mathbb Z$, an $h$-periodic insertion separates the two pole
ladders $\theta_s+nh$ ($n=0,\dots,3$), so the relation holds for each ladder separately. We
checked it in its original form, $L_s[\mu\,q^j_A\,\Op^Aq_i^{B[2]}]=0$, for all $30$ ordered cross
pairs, both ladders and all $(i,j)$: $540/540$.

\section{The $3\times3$ matrix}

\subsection{Rows on one ladder}
Three rows $(s,i,j)$ stacked give $N^{AB}$ with $N^{AB}\Delta=0$. So $\det N^{AB}=0$ for distinct
states, and $\det N^{AA}$ is the norm. Of the $816$ triples of rows, $732$ give $\det N^{AA}\neq0$
for all six states and the remaining $84$ give zero for all six. The simplest working triple uses
only the ladder of site~1. With
\begin{equation}
\br{f}=\oint_{C_1}\frac{du}{2\pi i}\,\frac{f(u)}{\prod_{\alpha}\prod_{n=0}^{3}(u-\theta_\alpha-nh)},
\end{equation}
where $C_1$ encircles $\theta_1+nh$, $n=0,\dots,3$, it reads
\begin{equation}
N^{AB}=\begin{pmatrix}
\br{q^1_A(u^2q_1^{B[-2]}-\nu_1\sh{-2}q_1^B)} & \br{q^1_A\,u\,q_1^{B[-2]}} & \br{q^1_A\,q_1^{B[-2]}}\\
\br{q^2_A(u^2q_1^{B[-2]}-\nu_1\sh{-2}q_1^B)} & \br{q^2_A\,u\,q_1^{B[-2]}} & \br{q^2_A\,q_1^{B[-2]}}\\
\br{q^1_A(u^2q_2^{B[-2]}-\nu_1\sh{-2}q_2^B)} & \br{q^1_A\,u\,q_2^{B[-2]}} & \br{q^1_A\,q_2^{B[-2]}}
\end{pmatrix}.
\end{equation}
With these rows, all six norms are non-zero and all $30$ cross determinants vanish. No cross matrix
vanishes identically.

\paragraph{The octets.} The two octets satisfy $T(8_1)-T(8_2)=(0,0,-\tfrac{129}{19})$, so the relation
forces the third column of $N^{8_18_2}$ to vanish. We find $\operatorname{rank}N^{8_18_2}=2$ and
\begin{equation}
\det N^{8_18_1}=\frac{53291113880694676031}{410708765061000000}.
\end{equation}
For $8_1$ the functions entering are $q_1=u-\tfrac{434}{285}$,
$q_2=u^3-\tfrac{41}{19}u^2+\tfrac{49084}{48735}$,
$q^1=u^5-\tfrac{305}{57}u^4+\tfrac{325295}{9747}u^2-\tfrac{151058891}{8333685}$ and
$q^2=u^3-\tfrac{84}{19}u^2+\tfrac{104791}{48735}$.

\subsection{Rows over the full contour}
Integrating every row over both ladders, $R_{i,j}$ vanishes identically exactly when $i<j$, as in
$\gl(4)$. Of the $84$ triples of the six remaining rows, $19$ give non-zero norms for all six
states at all three parameter points, and all $19$ give vanishing cross determinants.

\paragraph{Gauge.} The staggering gauge leaves the freedom $q_2\to q_2+a\,q_1$,
$q_3\to q_3+b\,q_1+c\,q_2$, $q^1\to q^1+d\,q^2+e\,q^3$ and $q^2\to q^2+f\,q^3$. Since $R_{i<j}=0$,
\begin{equation}
R_{21}\to R_{21}+a\,R_{11}+d\,R_{22},\qquad R_{32}\to R_{32}+c\,R_{22}+f\,R_{33},
\end{equation}
while $R_{11},R_{22},R_{33}$ are invariant. The only invariant triples with non-zero determinant are
therefore $\mathcal T_a=\{R_{11},R_{21},R_{22}\}$ and $\mathcal T_b=\{R_{22},R_{32},R_{33}\}$.
$\{R_{11},R_{22},R_{33}\}$ is also invariant, but its determinant vanishes. The two triples are
exchanged by $10\leftrightarrow\overline{10}$ together with a sign flip on the octets.

\section{Comparison with the Hilbert-space norm}

In the weight sector $\mu_A$ let $\Psi_R$ be the joint null vector of all coefficients of
$T_a-T_a^A$, and $\Psi_L$ the same for the transposed matrices. Every eigenline is
one-dimensional, and $\Psi_R$ is the highest-weight vector. In pivot gauge (first non-zero
component equal to one), $H_A=\Psi_L\cdot\Psi_R$. With $t=(\theta_1-\theta_2)/h$, computed
symbolically,
\begin{equation}
\widehat N_A=\frac{\det\mathcal T_a^{AA}}{\det\mathcal T_a^{\rm vac}}:\qquad
\widehat N_{10}=\widehat N_{1}=\frac{20}{3(49-25t^2)},\qquad
\widehat N_{\overline{10}}=\frac{160(t^2-1)}{3(25t^2-49)},
\end{equation}
whereas $H=1,2,2,8$ for $\mathrm{vac},10,\overline{10},1$, independent of $\theta_\alpha$ and $h$.
For the octets $H$ does depend on the parameters: $\tfrac{645}{128},\tfrac{86}{3}$ at $t=\tfrac9{19}$,
$\tfrac{405}{32},\tfrac{162}{25}$ at $t=-\tfrac7{11}$, and $\tfrac{830}{27},\tfrac{249}{50}$ at
$t=-\tfrac{33}{19}$ (with $h=\tfrac23$). The reason is that pivot gauge is not adapted to the octet
plane. None of the $19$ working full-contour triples gives a parameter-independent $\widehat N$.

\paragraph{$Q$-values at special points.} The $\theta_1\leftrightarrow\theta_2$ symmetry holds
exactly for the gauge-invariant functions $q_1$, $q_{12}=\mathrm{monic}[q_1q_2\sh{-2}-q_1\sh{-2}q_2]$,
the lowest dual $q^3$, and $q^{23}=\mathrm{monic}[q^2q^{3[2]}-q^{2[2]}q^3]$. It fails for the others.
For the non-octet states these four functions have only trivial roots ($\theta_\alpha$,
$\theta_\alpha+2h$) and midpoint roots $(\theta_1+\theta_2+kh)/2$. At points $\theta_\alpha+nh$
they therefore evaluate to products of $(t\pm k)/2$, and cannot produce the factors $7\pm5t$ of
$\widehat N$. Consistently, a bounded search of the ansatz
\begin{equation}
\widehat N_A\,\frac{\Pi_{\rm vac}}{\Pi_A}=c_A\,H_A,\qquad
\Pi=q^{a}(x_1)\,q^{bc}(x_2)\,q_1(x_3)\,q_{12}(x_4),
\end{equation}
over all dual labels, $x_i\in\{\theta_\alpha+nh:\ n=-2,\dots,4\}$ and both orientations, with
$c_A$ required independent of $t$, found no solution. The test was numerical: 30 digits, four
values of $t$, $1254\times168\times2$ candidates.

\section{Summary and open points}
\begin{enumerate}
\item The adjoint chain has three independent zero-twist differences of integrals of motion, and
the $3\times3$ functional determinant proves orthogonality for all states, same-weight octets
included.
\item The full-contour vacuum-normalised norm depends on the inhomogeneities and is not a
$Q$-dressed Hilbert norm in the tested class. A parameter-independent combination of rows, or rows
with $u^\beta$ insertions, remains to be found.
\item The relation to the twisted $6\times6$ norm is established in
\texttt{SU3\_Adjoint\_Twisted\_SoV\_L2}: the present $\mathcal T_a$ matrix is its upper-left minor.
\end{enumerate}

\end{document}
